% Part pages of Volume 1 (Chinese).  \partpage{N} prints part N; used by main_zh.tex and by the
% per-book entry files that build.py writes.  Generated once from the plan; edit freely.
\makeatletter
\@namedef{partpage@1}{\bookpart{language}{第一部分}{宏观的语言}{在进入历史之前，先准备好全书反复使用的语言：国民核算的恒等式、价格与利率、货币与银行的资产负债表、国际收支与汇率。这一部分不按历史顺序，只求让后面每一章的论证都有落脚点。}{%
\chref{ch:intro}\quad\nameref{ch:intro}\par
\chref{ch:accounts}\quad\nameref{ch:accounts}\par
\chref{ch:prices}\quad\nameref{ch:prices}\par
\chref{ch:money}\quad\nameref{ch:money}\par
\chref{ch:external}\quad\nameref{ch:external}}}
\@namedef{partpage@2}{\bookpart[公元前 3000 年—公元 600 年]{ancient}{第二部分}{古代世界}{货币、信用、税收与国家几乎同时出现。这一部分走过美索不达米亚与埃及、希腊与罗马、先秦与秦汉、古代印度与印度洋，追问最早的记账、铸币与物价管理怎样运作，以及后世的哪些宏观问题在那时已经出现。}{%
\chref{ch:mesopotamia}\quad\nameref{ch:mesopotamia}\par
\chref{ch:greekroman}\quad\nameref{ch:greekroman}\par
\chref{ch:qinhan}\quad\nameref{ch:qinhan}\par
\chref{ch:india}\quad\nameref{ch:india}}}
\@namedef{partpage@3}{\bookpart[600—1750 年]{earlymodern}{第三部分}{中世纪与近世}{从伊斯兰世界的汇票到宋代的纸币，从美洲白银到阿姆斯特丹与伦敦的国债，信用与货币在这一千多年里被一再重新发明。这一部分同样走进非洲、美洲与东亚，看不同的社会怎样解决同样的问题。}{%
\chref{ch:islamic}\quad\nameref{ch:islamic}\par
\chref{ch:medieval}\quad\nameref{ch:medieval}\par
\chref{ch:songming}\quad\nameref{ch:songming}\par
\chref{ch:silver}\quad\nameref{ch:silver}\par
\chref{ch:africa}\quad\nameref{ch:africa}\par
\chref{ch:finrevolution}\quad\nameref{ch:finrevolution}\par
\chref{ch:mercantilism}\quad\nameref{ch:mercantilism}\par
\chref{ch:tokugawa}\quad\nameref{ch:tokugawa}}}
\@namedef{partpage@4}{\bookpart[1750—1914 年]{classical}{第四部分}{古典世界}{工业革命让人均收入第一次持续增长，也催生了政治经济学。这一部分从斯密、休谟、李嘉图、马克思一直讲到维克塞尔与费雪，同时追问金本位时代的边缘国家与殖民地经历了什么。}{%
\chref{ch:smith}\quad\nameref{ch:smith}\par
\chref{ch:bullion}\quad\nameref{ch:bullion}\par
\chref{ch:say}\quad\nameref{ch:say}\par
\chref{ch:divergence}\quad\nameref{ch:divergence}\par
\chref{ch:marx}\quad\nameref{ch:marx}\par
\chref{ch:goldstandard}\quad\nameref{ch:goldstandard}\par
\chref{ch:periphery}\quad\nameref{ch:periphery}\par
\chref{ch:wicksell}\quad\nameref{ch:wicksell}}}
\@namedef{partpage@5}{\bookpart[1914—1945 年]{keynes}{第五部分}{战争、萧条与凯恩斯革命}{两次世界大战与大萧条动摇了古典的信条。这一部分从魏玛的恶性通胀、各国走出萧条的不同道路与苏联的计划经济讲起，再进入凯恩斯革命与战后国际货币秩序的设计。}{%
\chref{ch:weimar}\quad\nameref{ch:weimar}\par
\chref{ch:interwar}\quad\nameref{ch:interwar}\par
\chref{ch:depressionus}\quad\nameref{ch:depressionus}\par
\chref{ch:depressionworld}\quad\nameref{ch:depressionworld}\par
\chref{ch:alternatives}\quad\nameref{ch:alternatives}\par
\chref{ch:soviet}\quad\nameref{ch:soviet}\par
\chref{ch:generaltheory}\quad\nameref{ch:generaltheory}\par
\chref{ch:islm}\quad\nameref{ch:islm}\par
\chref{ch:warfinance}\quad\nameref{ch:warfinance}}}
\@namedef{partpage@6}{\bookpart[1945—1970 年]{goldenage}{第六部分}{战后黄金时代：东方与西方}{战后二十多年是西方的高增长期，也是计划经济与发展中国家各自探索道路的时期。这一部分并排讲述战后重建、布雷顿森林体系、菲利普斯曲线、增长与发展理论、东亚奇迹与成熟的计划经济。}{%
\chref{ch:reconstruction}\quad\nameref{ch:reconstruction}\par
\chref{ch:brettonwoods}\quad\nameref{ch:brettonwoods}\par
\chref{ch:mundell}\quad\nameref{ch:mundell}\par
\chref{ch:phillips}\quad\nameref{ch:phillips}\par
\chref{ch:consumption}\quad\nameref{ch:consumption}\par
\chref{ch:solow}\quad\nameref{ch:solow}\par
\chref{ch:development}\quad\nameref{ch:development}\par
\chref{ch:eastasia}\quad\nameref{ch:eastasia}\par
\chref{ch:planned}\quad\nameref{ch:planned}}}
\@namedef{partpage@7}{\bookpart[1968—1985 年]{stagflation}{第七部分}{滞胀与预期革命}{通胀与失业同时上升，让凯恩斯主义的政策菜单失灵。这一部分讲述自然率、货币主义、理性预期与可信度怎样重塑宏观经济学，也讲述浮动汇率、石油危机、拉美债务危机与恶性通胀。}{%
\chref{ch:naturalrate}\quad\nameref{ch:naturalrate}\par
\chref{ch:monetarism}\quad\nameref{ch:monetarism}\par
\chref{ch:floating}\quad\nameref{ch:floating}\par
\chref{ch:oil}\quad\nameref{ch:oil}\par
\chref{ch:ratex}\quad\nameref{ch:ratex}\par
\chref{ch:credibility}\quad\nameref{ch:credibility}\par
\chref{ch:budget}\quad\nameref{ch:budget}\par
\chref{ch:latamdebt}\quad\nameref{ch:latamdebt}\par
\chref{ch:hyperinflation}\quad\nameref{ch:hyperinflation}}}
\@namedef{partpage@8}{\bookpart[1982—2007 年]{moderation}{第八部分}{大缓和、新共识与转轨}{新凯恩斯框架与独立的央行带来了二十年的平静，计划经济则走向转轨。这一部分讲述新共识的形成、增长与失业理论、日本泡沫、欧元诞生、新兴市场危机，以及平静之下积累起来的风险。}{%
\chref{ch:rbc}\quad\nameref{ch:rbc}\par
\chref{ch:newkeynesian}\quad\nameref{ch:newkeynesian}\par
\chref{ch:consensus}\quad\nameref{ch:consensus}\par
\chref{ch:endogenous}\quad\nameref{ch:endogenous}\par
\chref{ch:labor}\quad\nameref{ch:labor}\par
\chref{ch:japanbubble}\quad\nameref{ch:japanbubble}\par
\chref{ch:transition}\quad\nameref{ch:transition}\par
\chref{ch:euro}\quad\nameref{ch:euro}\par
\chref{ch:emcrises}\quad\nameref{ch:emcrises}\par
\chref{ch:moderation}\quad\nameref{ch:moderation}\par
\chref{ch:minsky}\quad\nameref{ch:minsky}}}
\@namedef{partpage@9}{\bookpart[2008—2019 年]{aftercrisis}{第九部分}{金融危机之后}{2008 年的危机暴露了主流模型对金融的忽视。这一部分讲述危机本身、零利率下限与非常规政策、紧缩之争、欧债危机、长期停滞与异质性，以及中国的信贷驱动增长。}{%
\chref{ch:gfc}\quad\nameref{ch:gfc}\par
\chref{ch:frictions}\quad\nameref{ch:frictions}\par
\chref{ch:zlb}\quad\nameref{ch:zlb}\par
\chref{ch:abenomics}\quad\nameref{ch:abenomics}\par
\chref{ch:multipliers}\quad\nameref{ch:multipliers}\par
\chref{ch:eurocrisis}\quad\nameref{ch:eurocrisis}\par
\chref{ch:stagnation}\quad\nameref{ch:stagnation}\par
\chref{ch:heterogeneity}\quad\nameref{ch:heterogeneity}\par
\chref{ch:chinacredit}\quad\nameref{ch:chinacredit}}}
\@namedef{partpage@10}{\bookpart[2020—2026 年]{present}{第十部分}{疫情通胀与当下}{疫情、通胀回归与高债务把财政与货币的边界重新推到台前。这一部分讲述 2020 年以来的冲击与争论：财政主导、货币等级、数字货币、全球分工的碎片化，以及人口与 AI 带来的长期问题。}{%
\chref{ch:pandemic}\quad\nameref{ch:pandemic}\par
\chref{ch:inflation21}\quad\nameref{ch:inflation21}\par
\chref{ch:fiscaldominance}\quad\nameref{ch:fiscaldominance}\par
\chref{ch:hierarchy}\quad\nameref{ch:hierarchy}\par
\chref{ch:digitalmoney}\quad\nameref{ch:digitalmoney}\par
\chref{ch:fragmentation}\quad\nameref{ch:fragmentation}\par
\chref{ch:demography}\quad\nameref{ch:demography}}}
\@namedef{partpage@11}{\bookpart{synthesis}{第十一部分}{综合与展望}{最后一部分把全书的工具用于判断：怎样评估证据与数据，怎样从宏观走向资产价格，怎样复盘一个国家的长期历程，以及怎样有条件地谈论未来。}{%
\chref{ch:evidence}\quad\nameref{ch:evidence}\par
\chref{ch:datareliability}\quad\nameref{ch:datareliability}\par
\chref{ch:assetprices}\quad\nameref{ch:assetprices}\par
\chref{ch:longarcs}\quad\nameref{ch:longarcs}\par
\chref{ch:outlook}\quad\nameref{ch:outlook}}}
\makeatother
